\title{LongRoPE: Extending LLM Context Window Beyond 2 Million Tokens}

\begin{abstract}
Expanding the context window in large language models (LLMs) is highly sought after. Nevertheless, the achievable length of context windows remains constrained to roughly 128k tokens, primarily due to the significant expense of fine-tuning, the limited availability of extended textual data, and disruptive effects arising from unfamiliar token positions.
\end{abstract}

In this work, we present LongRoPE, a novel method that, for the first time, broadens the context window of pre-trained LLMs to an unprecedented \textbf{2048k} tokens. Remarkably, this is accomplished with as few as 1k fine-tuning steps at training lengths up to 256k, all while preserving the models’ short context window capabilities. The approach relies on three primary contributions: (i) we identify and leverage two distinct types of non-uniformities present in positional interpolation by performing an efficient search, which yields superior fine-tuning initialization and allows for an 8$\times$ window expansion even without fine-tuning; (ii) a progressive extension framework is adopted in which the model is first fine-tuned at a 256k context length and subsequently subjected to an additional positional interpolation, extending capacity to 2048k tokens; (iii) the LongRoPE method undergoes adjustment at the 8k length to restore its short context window performance. Comprehensive experimental evaluation using LLaMA2 and Mistral, over multiple tasks, confirms the robustness of this method. The modifications required by LongRoPE involve only minor changes to the positional embeddings, leaving the original model architecture and its compatibility with prior optimizations largely intact. The code will be made available at \texttt{https://github.com/microsoft/LongRoPE}
\end{abstract}

\section{Introduction}

Large Language Models (LLMs) have achieved significant progress across a range of tasks~\cite{gpt4,llama2}; however, their effectiveness is constrained by the size of their context window, such as LLaMA2's 4096-token capacity~\cite{llama2}. When inputs exceed this limit, model performance degrades because the network has not encountered such positional contexts during training. This limitation creates obstacles for use cases that demand processing large amounts of context, such as in-context learning involving many examples~\cite{Huang2023FewerIM} and applications like LLM agents~\cite{park2023generative,madaan2023selfrefine}.

Recent studies have demonstrated that the context window of pre-trained LLMs can be expanded up to approximately 128k by fine-tuning the models on longer sequences~\cite{longlora,pi,yarn,zhang2024soaring,liu2023scaling}. However, three primary challenges hinder further increases in context length. First, introducing new position indices that have not been previously trained results in substantial numbers of problematic values, causing out-of-distribution difficulties and complicating the fine-tuning process~\cite{pi}. This obstacle becomes pronounced when the context window is extended from 4k to $>$1000k, where over 90\% of token positions are new. Second, effective fine-tuning typically necessitates texts that match the expanded context lengths. Unfortunately, datasets containing sequences longer than 1000k are scarce, and training on such exceptionally lengthy sequences is computationally intensive, demanding considerable training time and significant GPU resources. Third, upon expanding the context to extremely large windows, the attention mechanism becomes diluted as it must distribute focus across an enormous number of token positions, resulting in reduced performance on tasks involving the original short contexts~\cite{pi}.

A commonly used strategy to address the initial challenge involves interpolating RoPE positional embeddings~\cite{rope,pi}, which effectively rescales new positional indices to align with the original pretrained interval, as depicted in Fig.\ref{fig:overview}. Position Interpolation (PI)~\cite{pi} performs a linear interpolation of RoPE’s rotary angles according to the expansion ratio. In contrast, NTK~\cite{ntk,dynamicntk} introduces unequal methods for interpolating and extrapolating among different RoPE dimensions. YaRN~\cite{yarn} takes a frequency-oriented perspective, dividing RoPE dimensions into three distinct groups and assigning extrapolation, NTK, or linear interpolation strategies to each. Nonetheless, the positional embedding mechanism within Transformer architectures contains \emph{complex and uneven distributions of information entropy}. Contemporary techniques do not sufficiently exploit this intricate non-uniformity, resulting in a degree of information loss and imposing constraints on achievable context window size.

Section~\ref{sec:analysis} empirically demonstrates two principal results: \textbf{(1)} Successful positional interpolation necessitates accounting for two distinct types of non-uniformity—those arising from differences across RoPE dimensions and those associated with specific token locations. Reduced interpolation proves beneficial for tokens at the beginning and lower RoPE dimensions, though optimal strategies vary based on the intended extended sequence length.  
\textbf{(2)} Incorporating these non-uniformities into the interpolation process enables better preservation of the information encoded by the original RoPE, especially regarding salient dimensions and token placements. This approach reduces information loss during interpolation and thus yields superior initialization when fine-tuning models. Additionally, it supports up to an 8$\times$ increase in sequence length without the need for fine-tuning.

Building upon these insights, we propose \textbf{LongRoPE}, a robust approach designed to increase the LLM context window to over 2 \emph{million} tokens. The framework of LongRoPE is founded on three main advancements. To begin with, LongRoPE leverages multidimensional positional interpolation with non-uniform characteristics. It discovers optimal rescale factors for RoPE's rotational angles, assigning these for each dimension according to the respective token positions. Since the process of determining rescale factors is subject to an exponentially enlarging search space as the extension ratio rises, LongRoPE employs an evolutionary search strategy enhanced by two dedicated optimization methods, thereby improving search performance. Figure~\ref{fig:overview} illustrates the resulting rescaled RoPE, obtained through this procedure.

Subsequently, {LongRoPE} adopts a resource-efficient, stepwise extension approach to reach a 2048k context window, avoiding the requirement for direct fine-tuning on ultra-long sequences, which are rare and difficult to procure. The procedure starts by identifying the feasible maximum length of 256k for the pre-trained LLM, followed by fine-tuning at this length. Owing to our non-uniform positional interpolation method, which supports up to an 8$\times$ expansion even without further fine-tuning, we then perform a secondary search to determine new RoPE rescaling parameters on the already fine-tuned, longer-context LLM. This process enables {LongRoPE} to support a 2048k context window for both LLaMA2 and Mistral~\cite{mistral}.

In order to address potential losses in performance when dealing with shorter context windows, {LongRoPE} further refines the RoPE rescale parameters within the expanded LLM. Using an approach analogous to increasing the context length from 256k to 2048k, we instead decrease it to 4k and 8k windows for the LLM fine-tuned at 256k, applying our search algorithm to minimize positional interpolation. During inference, when the sequence length falls below 8k, RoPE is updated utilizing the identified rescale factors.

A comprehensive set of experiments conducted on multiple LLMs and a range of long-context tasks highlights the strong performance of our approach. The results indicate that LongRoPE sustains low perplexity across evaluation lengths spanning from 4k up to 2048k, exceeds 90\% accuracy in passkey retrieval tasks, and matches the accuracy of existing models on benchmarks limited to a 4096 context window. Additionally, LongRoPE is compatible with any LLM that utilizes RoPE embeddings. We plan to make both our codebase and the LongRoPE-2048k models publicly available.

\section{Non-uniformity in Positional Interpolation}
\label{sec:analysis}

\subsection{Preliminary}

Transformer models require explicit positional information, often in the form of position embedding, to represent the order of input tokens. Our work focuses on the RoPE~\cite{rope} position embedding, which is widely used in recent  LLMs. For a token at position index $n$, its corresponding RoPE encoding can be simplified as follows:

\begin{equation}
		\fontsize{7.8}{7.8} \selectfont
	\label{eq:rope}
	\left[ cos(n\theta_0), sin(n\theta_0), cos(n\theta_1), \cdot\cdot\cdot, cos(n\theta_{d/2-1}), sin(n\theta_{d/2-1}) \right]   
\end{equation}

Here, $d$ denotes the size of the embedding vector, and $n\theta_i$ specifies the rotary phase associated with the token found at position $n$. The parameter $\theta_i=\theta^{-2i/d}$ characterizes the rotation frequencies utilized. Typically, RoPE employs a base value of $\theta$ set to 10000.

\textbf{\emph{Context window extension ratio $s$} and positional interpolation}. The ratio $s$ is introduced as a measure, defined by the quotient of the expanded context length $L'$ over the original context length $L$, i.e., $s=\frac{L'}{L}$.

To extend the context window from $L$ to $L'$, current positional interpolation methods suggest downscaling rotation frequencies $\theta_i$ based on extension ratio $s$. Let $\beta=\theta^{2/d}$, and $\lambda$ denote the actual rescale factor related to $s$, we   unify these positional interpolation methods as follows:

\begin{equation}	
	\fontsize{7.5}{7.5} \selectfont
	\label{eq:unified_rope}
	\begin{aligned}
		\left[cos\left(\frac{n}{\lambda(\beta)^0}\right), sin \left(\frac{n}{\lambda(\beta)^0}\right), cos\left(\frac{n}{\lambda(\beta)^1}\right), \cdot\cdot\cdot,  sin \left(\frac{n}{\lambda(\beta)^{d/2-1}}\right) \right]   
	\end{aligned}
\end{equation}

\textbf{Linear positional interpolation (PI)}. PI~\cite{pi} proposes to perform linear interpolation of position indices, constrained within the original pre-training sequence length. When applying an extension ratio $s$, the rotation angles assigned to each position are uniformly scaled down by $\lambda=s$ in every dimension of RoPE. This approach, however, compresses positional information, leading to difficulty in differentiating tokens that are nearby in sequence. Consequently, PI is generally less effective when the extension ratio is high.

\textbf{NTK-based interpolation and extrapolation}. In their work, ~\cite{ntk,dynamicntk} re-examine RoPE from the standpoint of information encoding, employing Neural Tangent Kernel (NTK) theory~\cite{jacot2018neural,tancik2020fourier}. Addressing the challenge of crowded positional encodings inherent to PI, these authors propose to evenly spread interpolation constraints over the dimensions of RoPE. This technique involves reducing the scaling of lower (high frequency) dimensions and increasing it for higher (low frequency) dimensions, thereby facilitating both interpolation and extrapolation of position—specifically using $\lambda=s^{i}$. Building on this, the enhanced dynamic NTK approach~\cite{dynamicntk} modulates the extension ratio at each sequence position depending on the actual sequence length. In contrast to PI, which relies on model fine-tuning, NTK-sensitive strategies are capable of enlarging context windows without additional fine-tuning, although they generally support up to a 4$\times$ extension ratio.

\textbf{YaRN}~\cite{yarn} offers a notable enhancement in positional interpolation for RoPE by categorizing its dimensions into three separate groups based on their frequencies, each utilizing distinct interpolation techniques. In this method, dimensions associated with high frequencies are subject to extrapolation ($\lambda$=1), low frequency dimensions apply linear interpolation (PI), and intermediate frequencies leverage NTK. A central feature of YaRN is the way it clusters RoPE dimensions; however, this grouping process is currently driven by manual empirical observations, which may yield less-than-ideal outcomes for novel LLM architectures.

\subsection{Study on Non-uniform Positional Interpolation}

Drawing motivation from NTK and YaRN, we observe their improvements stemming from the introduction of non-linear elements, particularly the treatment of distinct frequencies within RoPE dimensions to enable tailored interpolation and extrapolation. Nevertheless, existing non-linear functions are predominantly based on manually crafted heuristics. This prompts us to ask: (1) Does the current approach to positional interpolation represent the best possible method? (2) Might there be non-linear strategies that have not yet been investigated?

To address these issues, we apply evolution search (refer to Sec.~\ref{sec:method}) to identify improved non-uniform positional interpolation schemes for LLaMA2-7B. This process leverages perplexity as the optimization metric and utilizes 5 randomly selected samples from the PG19~\cite{pg19} validation set. Our experiments yield the following principal observations.

\textbf{Finding 1}: \textit{RoPE dimensions exhibit substantial non-uniformities, which are not effectively handled by current positional interpolation methods.}

We perform a search to determine the optimal value of $\lambda$ for every RoPE dimension as specified in Eq.~\ref{eq:unified_rope}. Table~\ref{tbl:exp1} presents a perplexity comparison for LLaMA2-7B across various approaches on the PG19 and Proof-pile~\cite{proof-pile} test datasets, evaluated without any fine-tuning. Our empirically derived settings yield notable perplexity reductions, indicating that both the standard linear (PI) and the more complex non-uniform approaches (Dynamic-NTK and YaRN) are not fully optimized. Interestingly, YaRN achieves worse results than PI and NTK on the PG19 set, because it does not extend to the desired context window in the absence of fine-tuning for LLMs. To illustrate, with an 8k context window, the perplexity using YaRN rises steeply beyond the 7k token position.

Our investigation reveals that the rescaling coefficients $\lambda$ in Eq.~\ref{eq:unified_rope} are not uniform across positions, in contrast to the constant scaling factor $s$ used in PI, the deterministic computation of NTK, and the segmented scaling applied in YaRN. Employing these position-dependent rescale factors yields a substantial enhancement in LLaMA2's language modeling metrics (measured via perplexity) when utilizing extended context windows of 8k and 16k, without necessitating additional fine-tuning. This improvement is attributed to the modified positional encoding's ability to retain the essential characteristics of the original RoPE, with particular emphasis on critical dimensions, thereby alleviating the model's challenge in differentiating tokens situated in proximate positions.

\textbf{Finding 2}: \textit{RoPE for the initial tokens in the input sequence should be extrapolated with less interpolation.} 

We propose that the RoPE should interpolate less for the first $\hat{n}$ tokens in input sequences. This assumption is based on their tendency to receive high attention scores, thus playing a significant role in attention mechanisms, as reported in Streaming LLM~\cite{streamingllm} and LM-Infinite~\cite{han2023lminfinite}. To assess this, we enlarged the context window to 8k and 16k utilizing PI and NTK, ensuring that the initial $\hat n$ (0, 2, ..., 256) tokens are preserved without interpolation. When $\hat n$ equals zero, we use the standard PI and NTK approaches. Table~\ref{tbl:tokenposition} reveals two primary findings: \textbf{(1)} omitting position interpolation for the initial tokens enhances performance; \textbf{(2)} the best value of $\hat n$ varies according to the desired extension length.

\textbf{Finding 3}: \textit{Non-uniform positional interpolation effectively extends LLM context window in both fine-tuning and non-fine-tuning settings.} 

Although our explorations demonstrate that applying our optimized non-uniform position interpolation leads to notable gains in extension capability at 8k and 16k context lengths—even in the absence of fine-tuning—extending to even longer sequences necessitates model fine-tuning. To address this, we conduct fine-tuning of LLaMA2-7B using our discovered RoPE method, scaling up to a 64k context length (refer to the Appendix for experimental configurations). Table~\ref{tbl:finetune} highlights that our approach offers clear advantages over both PI and YaRN, outperforming them pre- and post-fine-tuning of LLaMA2-7B. The observed improvements are attributable to our approach's capacity for non-uniform positional interpolation, which more effectively preserves information and furnishes superior initialization for the fine-tuning process.

\textbf{Summary}. This work identifies two forms of non-uniformity—variations in RoPE dimensions and in token positions—and shows that harnessing these non-uniformities within positional interpolation strategies leads to substantial improvements in context extension for large language models.

\section{LongRoPE}

\label{sec:method}
Inspired by these results, we propose LongRoPE, a method that begins by implementing a novel search algorithm designed to leverage both forms of non-uniformity. Subsequently, this approach enables the extension of LLM context windows to surpass 2 million tokens.

\subsection{Problem Formulation}

These two types of non-uniformity greatly expand the range of possible solutions and complicate the optimization process. To manage this challenge, we reformulate the multidimensional non-uniform position interpolation optimization as a search problem.

For a LLM targeting a  context window size of $L'$ and lengthy input documents $\mathbf{X}$, where each $\mathbf{x}\in\mathbf{X}$ surpasses $L'$ in token length, we denote the original rotary angle of the $i^{th}$ dimension in RoPE embedding at token position $n$ as  $\frac{n}{\beta^i}$. The optimization problem is then formulated as follows:

\begin{equation}
\begin{gathered}
		\label{eq:problem}

\underset {\mathbf{x}\in\mathbf{X}; \, |\mathbf{x}|\ge L'}{ \operatorname {arg\,min} } \,  \mathcal{L}\left(\text{LLM}(\text{RoPE}, \mathbf{X})\right), \text{ with} 
\\
\scalebox{0.93}{
$\underset {\begin{subarray}{l} i=0,\cdots,\frac{d}{2}-1;\\ n\in[ 0,|\mathbf{x}|);\end{subarray}}{\text{RoPE}_ {(n)}}= {\left[\cdots, \cos\left(\mathbb{I}(\hat{\lambda}_{i}, \hat{n}) \cdot \frac{n}{\beta^i}\right), \sin\left(\mathbb{I}(\hat{\lambda}_{i}, \hat{n}) \cdot \frac{n}{\beta^i}\right), \cdots \right] }$}\\
	\text{where} \quad \mathbb{I}(\hat{\lambda}_{i},\hat{n})=\begin{cases}
	1&\text{if } n<\hat{n}\\
{\frac{1}{\lambda}_{i}}&\text{if } n\geq\hat{n}
\end{cases}
	\end{gathered}
\end{equation}

In this approach, we define a set of rescaling factors, $\mathbb{I}(\hat{\lambda}_{i},\hat{n})$, to account for the two distinct types of non-uniformities present. Here, $\hat{\lambda_i}$ corresponds to the non-uniformity in RoPE dimensions, while $\hat{n}$ refers to non-uniformities in token positions. The factor $\mathbb{I}(\hat{\lambda}_{i},\hat{n})$ is employed to adjust the rotation angle for each $i^{th}$ RoPE dimension: $\hat{\lambda}_{i}$ serves as the scaling coefficient and $\hat{n}$ determines the token position threshold. For token positions less than $\hat{n}$, the rescale factor $\hat{\lambda}_{i}$ is not applied, so the standard RoPE rotary angle $\frac{n}{\beta^i}$ remains. Once the token position satisfies $n \ge \hat{n}$, the rescale factor is then introduced.

With a specified target context window size $L'$, our aim is to determine the set of optimal rescale factors ($\mathbb{I}(\hat{\lambda}_{0},\hat{n})$, $\mathbb{I}(\hat{\lambda}_{1},\hat{n})$, ..., $\mathbb{I}(\hat{\lambda}_{i},\hat{n})$, ...) across the $1^{st}$ through $d^{th}$ dimensions of RoPE. This adjustment is intended to ensure that the resulting $\text{LLM}$, incorporating the modified $\text{RoPE}$, yields the lowest possible next token prediction loss, denoted by $\mathcal{L}$ (namely, perplexity), on input sequences $\mathbf{X}$ of length $L'$.

\subsection{Searching the Non-uniform Position Interpolation}

To address the issue posed in Eq.~\ref{eq:problem}, we propose a straightforward but powerful technique that identifies the ideal RoPE rescaling parameters, thereby maximizing the utilization of multidimensional irregularities present in position embedding.

\textbf{Search space}. We establish an extensive search space that accommodates both identified non-uniformities. Table~\ref{tbl:searchspace} provides a summary of this search space. In particular, our approach permits a unique rescale factor to be determined for each dimension of the RoPE embedding. To simplify the formulation of the search space, we perform a direct search over $\lambda_i$ and $\hat{n}$ rather than over $\mathbb{I}(\hat{\lambda}_{i},\hat{n})$, using the relationship $\hat{\lambda}_{i}=1/\lambda_i$. As detailed in Table~\ref{tbl:searchspace}, the value of $\lambda_i$ spans from a lower bound of 1.0 (which corresponds to plain extrapolation) up to an upper limit of $s\times1.25$ (offering more generous interpolation compared to PI), with increments of 0.01. Here, $s$ denotes the intended expansion ratio for the context window.

The parameter $\hat{n}$ determines how many of the starting token positions are preserved using the standard RoPE embedding, rather than employing position interpolation. In practice, we allow $\hat{n}$ to take values from the set \{0, 1, 2, 4, 8, 12, 16, 20, 24, 28, 32, 64, 128, 256\} during the search process. If $\hat{n}=0$, then every token position relies exclusively on the searched rescale factors.

\textbf{Evolution-based search}. The search space outlined in Table~\ref{tbl:searchspace} includes a vast array of positional interpolation strategies, making it extremely challenging to explore efficiently. For instance, an extension ratio of $s=4\times$ yields $400^{128/2}\times14=4\times10^{167}$ possible configurations. As the extension ratio increases, the number of options grows exponentially. To overcome these obstacles, we employ evolution search~\cite{guo2020single} and incorporate two optimization strategies to significantly enhance the efficiency of the search process. The comprehensive search workflow is described in Algorithm~\ref{alg:search}.

\textit{Enhanced strategy for generating the initial population}. Rather than selecting $P$ rescale factors purely at random for initialization, we explicitly include the RoPE rescale factors derived from PI, NTK, and YaRN within the first population. To fill the rest of the population ($P$-3), we introduce random mutations to these three rescale factors, applying mutations with a probability $p$.

\textit{Monotonically non-decreasing constraint}. Upon establishing the initial population, we evaluate the LLM perplexity for every individual. For this, the relevant RoPE rescale factors are applied to the target LLM, and the perplexity of input $\mathbf{X}$ is measured. The highest-scoring $k$ individuals are then chosen to act as parents in the evolutionary process. Nonetheless, the immense search space presents a challenge: straightforward mutation and crossover can result in unproductive exploration of suboptimal candidates, which in turn leads to superfluous perplexity evaluations. This problem is exacerbated for large $L'$, where each perplexity computation is highly resource-intensive.

In order to mitigate this issue, we enforce a monotonicity constraint such that the sampled RoPE rescale coefficients are non-decreasing: $\lambda_i \leq \lambda_{i+1}$. Only those RoPE configurations meeting this requirement are utilized in the LLM perplexity assessment, thereby substantially lowering the associated computational overhead. More precisely, we stipulate that $\lambda_i$ exhibits a monotonic increase across the RoPE dimensions (with $i$ ranging from 0 to 63). This monotonicity in the dimension index is motivated by the principles of NTK theory~\cite{jacot2018neural,tancik2020fourier,ntk}, which indicate that lower-dimensional indices, corresponding to higher frequencies, necessitate less interpolation (implying a smaller value for $\lambda_i$), whereas higher-dimensional indices, associated with lower frequencies, can accommodate greater interpolation (which entails a larger $\lambda_i$).

\begin{algorithm}[t]
	\small
	\caption{The search algorithm}
	\textbf{Input:} target LLM, input samples $\mathbf{X}$, population size $P$, mutation size $N_1$, crossover size $N_2$, max iterations $\mathcal{T}$, mutation probability $p$\\

	\begin{algorithmic}[1]
		\label{alg:search}
		\STATE Initialize Top-k as empty set;	
		\STATE $\text{P}_0$ = \textit{Initialize\_population\_with\_optimization}($P$, $\mathbf{X}$, $p$); 
		\FOR{$i$ from $1$ to $\mathcal{T}$}
		\STATE \textit{Compute\_perplexity}(LLM, $\text{P}_{i-1}$, $\mathbf{X}$);
		\STATE  Top-k $\leftarrow$ \textit{Update\_Topk}(Top-k, $\text{P}_{i-1}$);
		\STATE $\text{P}_{mutation}$ = \textit{Mutation\_with\_mono\_constraint}(Top-k, $N_1$, $p$); 
		\STATE $\text{P}_{crossover}$ = \textit{Crossover\_with\_mono\_constraint}(Top-k, $N_2$);
		\STATE $\text{P}_i$ = $\text{P}_{mutation}$ $\cup$ $\text{P}_{crossover}$ $\cup$ Top-k;
		\ENDFOR
		\STATE Return the member of Top-k exhibiting the minimum perplexity;
	\end{algorithmic}
\end{algorithm}

\textbf{8$\times$ extension without fine-tuning}. Through evolutionary search, our approach successfully discovers non-uniform RoPE rescaling parameters, selectively retaining crucial dimensions and positional information to reduce information loss caused by interpolation. As shown in Fig.\ref{fig:non-ft}, our technique enables the context window of LLaMA2 to be expanded from 4k to 32k without the need for fine-tuning. By comparison, alternative strategies like PI as well as non-uniform NTK and YaRN, lead to a sharp increase in perplexity beyond a 2$\times$ expansion.

\subsection{Extending LLM Context Window to 2048K}
\label{sec:extendto2048k}

\textbf{Progressive extension to 2048k}. In this section, we present our approach for scaling the context window of pre-trained LLMs beyond the standard 4k tokens, reaching up to 2048k tokens. Our non-uniform positional interpolation enables us to achieve an 8$\times$ context window increase without necessitating fine-tuning. However, for substantially larger expansions (such as 512$\times$), fine-tuning becomes essential. One strategy entails identifying suitable RoPE rescaling factors tailored to the 2048k context length, followed by fine-tuning. Nonetheless, this strategy is hindered by the immense computational costs involved. Furthermore, our empirical observations indicate that effective fine-tuning with such large extension ratios is quite difficult (refer to Appendix).

Fortunately, LongRoPE demonstrates efficacy with both base and fine-tuned extended LLMs. As such, we propose a progressive, resource-efficient approach that reaches the desired 2048k target in only 1,000 fine-tuning steps, each conducted at a maximum sequence length of 256k.

$\Diamond$ \textit{LongRoPE search for scaling pre-trained LLM to 256k}. Using LLaMA2 for illustration, we implement a search process aiming to reach context window sizes of 128k and 256k. At this phase, the enlargement rates correspond to 32$\times$ and 64$\times$, accordingly.

$\Diamond$ \textit{Fine-tuning to 256k}. Subsequently, we adapt the pre-trained LLM to reach a context window size of 256k through a two-phase fine-tuning procedure. Initially, LLaMA2 undergoes 400 fine-tuning steps with RoPE rescaling configured for 128k. Upon completing this phase, we update the RoPE rescale factors to 256k in the resulting checkpoint, followed by a further 600 fine-tuning steps. Compared to immediate fine-tuning towards 256k, this staged approach demonstrates improved efficiency.

$\Diamond$ \textit{Expanding fine-tuned extended LLM to 2048k using LongRoPE search}. In the final step, a subsequent search is applied to the LLM that has been fine-tuned for a 256k-length context. This leads to a context window of 2048k, achieved without additional fine-tuning. The resulting context extension represents a $512\times$ increase.

\textbf{Recovering performance within the shorter context window}. Upon expanding the context window to a substantial length of 2048k, we observe that model accuracy diminishes within the scope of the original context window. This challenge is well-documented in the context of positional interpolation~\cite{pi}: compressing the positional embeddings of the initial context window into a reduced subspace constricts their expressivity, thus degrading overall language modeling capabilities. Specifically, at an extension factor of 512$\times$, the positions confined to the initial 4k window are subject to pronounced overcrowding.

To address this issue, an additional evolution search is conducted on the extended LLM in order to refine RoPE rescale parameters specifically for shorter context lengths (such as 4k and 8k). The upper bound for $\lambda$ in the search is decreased, since shorter sequences demand less positional interpolation. At inference time, the LLM modifies the relevant RoPE rescale factors on-the-fly.

\section{LongRoPE}

\label{sec:method}
Building upon these insights, we propose LongRoPE. This approach begins with the implementation of a streamlined search algorithm that leverages both sources of non-uniformity. Utilizing this algorithm, LongRoPE expands the context window for LLMs, pushing it past the 2 million token threshold.

\subsection{Problem Formulation}

These two types of non-uniformity result in an expansive solution space and complicate the optimization process. To mitigate these challenges, we reformulate the optimization of multidimensional non-uniform position interpolation as a search problem.

For a LLM targeting a  context window size of $L'$ and lengthy input documents $\mathbf{X}$, where each $\mathbf{x}\in\mathbf{X}$ surpasses $L'$ in token length, we denote the original rotary angle of the $i^{th}$ dimension in RoPE embedding at token position $n$ as  $\frac{n}{\beta^i}$. The optimization problem is then formulated as follows:

\begin{equation}
\begin{gathered}
		\label{eq:problem}

\underset {\mathbf{x}\in\mathbf{X}; \, |\mathbf{x}|\ge L'}{ \operatorname {arg\,min} } \,  \mathcal{L}\, \left( \text{LLM} (\text{RoPE}, \mathbf{X})\right),	\text{with} 
\\
\scalebox{0.93}{
$\underset {\begin{subarray}{l} i=0,\cdot\cdot,\frac{d}{2}-1;\\ n\in[ 0,|\mathbf{x}|);\end{subarray}}{\text{RoPE}_(n)}= {\left[\cdots, \cos\left(\mathbb{I}(\hat{\lambda}_{i}, \hat{n})\frac{n}{\beta^i}\right),  \sin\left(\mathbb{I}(\hat{\lambda}_{i}, \hat{n})\frac{n}{\beta^i}\right),\cdots\right] }$}\\
	\text{where} \quad \mathbb{I}(\hat{\lambda}_{i},\hat{n})=\begin{cases}
	1 &\text{ if } n<\hat{n} \\
{\frac{1}{\lambda}_{i}} &\text{ if } n\geq\hat{n}
\end{cases}
	\end{gathered}
\end{equation}

In this approach, we define a set of scaling factors, $\mathbb{I}(\hat{\lambda}_{i},\hat{n})$, to address two distinct types of non-uniformities. Here, $\hat{\lambda_i}$ represents the degree of non-uniformity across RoPE dimensions, while $\hat{n}$ refers to variation over token positions. The function $\mathbb{I}(\hat{\lambda}_{i},\hat{n})$ specifically modifies the rotation angle associated with the $i^{th}$ RoPE dimension: $\hat\lambda_{i}$ is utilized as the scaling factor, and $\hat{n}$ serves as the threshold for token positions. For the first $\hat{n}-1$ tokens, the scaling factor $\hat\lambda_{i}$ is omitted, preserving the standard RoPE rotary angle of $\frac{n}{\beta^i}$. When token positions satisfy $n\ge\hat{n}$, the scaling factor is engaged.

For a specified context window length $L'$, our goal is to determine the most suitable rescale factors, denoted as ($\mathbb{I}(\hat{\lambda}_{0},\hat{n})$, $\mathbb{I}(\hat{\lambda}_{1},\hat{n})$, ... $\mathbb{I}(\hat{\lambda}_{i},\hat{n})$, ...), across all $d$ dimensions of the RoPE. By doing so, we aim to enable the target $\text{LLM}$ utilizing the rescaled $\text{RoPE}$ to attain the lowest possible next-token prediction loss, $\mathcal{L}$ (equivalently, minimizing perplexity), when processing input sequences $\mathbf{X}$ whose token count is $L'$.

\subsection{Searching the Non-uniform Position Interpolation}

In order to address the issue presented in Eq.~\ref{eq:problem}, we present a straightforward but powerful approach that identifies ideal RoPE rescale factors, thereby harnessing the complex, multidimensional variations inherent in position embedding.

\textbf{Search space}. We construct an extensive search space that incorporates both forms of non-uniformity. This space is summarized in Table~\ref{tbl:searchspace}. In particular, for the RoPE embedding, we permit the optimization of distinct rescale factors for each dimension. To simplify the structure of the search space, we opt to directly search over $\lambda_i$ and $\hat{n}$, rather than exploring $\mathbb{I}(\hat{\lambda}_{i},\hat{n})$, where $\hat{\lambda}_{i}=1/\lambda_i$. Table~\ref{tbl:searchspace} demonstrates that $\lambda_i$ is selectable from a lower bound of 1.0 (which corresponds to pure extrapolation) up to an upper bound of $s\times1.25$ (enabling more pronounced interpolation than PI), incremented by steps of 0.01. Here, $s$ represents the intended ratio for expanding the context window.

The parameter $\hat{n}$ determines how many of the starting token positions keep the standard RoPE embedding without any position interpolation applied. In our experiments, we permit $\hat{n}$ to take on values from the set \{0, 1, 2, 4, 8, 12, 16, 20, 24, 28, 32, 64, 128, 256\}. If $\hat{n}=0$, this implies that position interpolation—with learned rescale factors—is used for every token position.

\textbf{Evolution-based search}. The search space detailed in Table~\ref{tbl:searchspace} encompasses a vast array of positional interpolation configurations, making it inherently difficult to explore efficiently. For instance, employing an extension factor of $s=4\times$ results in $400^{128/2}\times14=4\times10^{167}$ possible options. As the extension ratio increases, the search space grows exponentially. To mitigate this complexity, we employ evolution search~\cite{guo2020single} along with two key optimization strategies that significantly enhance search efficiency. The complete search process is presented in Algorithm~\ref{alg:search}.

\textit{Optimized initial population generation}. Rather than selecting $P$ rescale factors purely at random to form the initial population, we explicitly include the three RoPE rescale factors associated with PI, NTK, and YaRN among the initial individuals. For the other $P$-3 members of the population, we generate them by applying random mutations to these three rescale factors with mutation probability $p$.

\textit{Monotonically non-decreasing constraint}. Once the initial population is established, we evaluate the LLM perplexity of every candidate. This involves applying the relevant RoPE rescale factors within the target LLM and determining the perplexity on input $\mathbf{X}$. The top-k performing candidates are selected as parents for subsequent evolutionary steps. Despite this, the extensive search landscape means that straightforward application of mutation and crossover may direct the process toward suboptimal regions, resulting in many fruitless perplexity evaluations. This inefficiency becomes pronounced as $L'$ grows, given that each perplexity computation demands significant inference time.

To tackle this challenge, a monotonicity constraint is enforced on the sampled RoPE rescaled factors such that $\lambda_i \le \lambda_{i+1}$. Only those RoPE configurations fulfilling this criterion are utilized in the LLM during perplexity testing, thereby substantially curbing the computational expense involved in searching. Concretely, $\lambda_i$ must show a monotonic increase across the RoPE dimensions ($i=0,\ldots,63$). The rationale for this dimension-wise monotonicity comes from NTK theory~\cite{jacot2018neural,tancik2020fourier,ntk}, which posits that dimensions associated with higher frequencies (lower $i$) require minimal interpolation—implying smaller $\lambda_i$ values—while dimensions with lower frequencies (higher $i$) accommodate more interpolation, corresponding to larger $\lambda_i$.

\begin{algorithm}[t]
	\small
	\caption{The search algorithm}
	\textbf{Input:} target LLM, input samples $\mathbf{X}$,   population size $P$, mutation size $N_1$, crossover size $N_2$, max iterations $\mathcal{T}$, mutate probability $p$\\

	\begin{algorithmic}[1]
		\label{alg:search}
		\STATE Initialize Top-k to $\phi$;
		\STATE Set initial population $\text{P}_0$ via \textit{Initialize\_population\_with\_optimization} ($P$, $\mathbf{X}$, $p$); 
		\FOR{$i$=1 to $\mathcal{T}$}
		\STATE Execute \textit{Compute\_perplexity} (LLM, $\text{P}_{i-1}$, $\mathbf{X}$);
		\STATE Refresh Top-k using \textit{Update\_Topk} (Top-k, $\text{P}_{i-1}$);
		\STATE Obtain $\text{P}_{mutation}$ by running \textit{Mutation\_with\_mono\_constraint} (Top-k, $N_1$, $p$); 
		\STATE Derive $\text{P}_{crossover}$ via \textit{Crossover\_with\_mono\_constraint} (Top-k, $N_2$);
		\STATE Formulate $\text{P}_i$ by merging $\text{P}_{mutation}$, $\text{P}_{crossover}$, and Top-k;
		\ENDFOR
		\STATE Output the member in Top-k with minimum perplexity;
	\end{algorithmic}
\end{algorithm}

\textbf{8$\times$ extension without fine-tuning}. Utilizing an evolutionary search, our approach discovers optimal non-uniform RoPE rescale factors, retaining critical dimensional and positional attributes to reduce the loss of information caused by interpolation. As illustrated in Fig.\ref{fig:non-ft}, this strategy successfully broadens LLaMA2’s context window from 4k to 32k with no need for fine-tuning. Conversely, other techniques like PI, as well as non-uniform NTK and YaRN, result in a significant increase in perplexity beyond a 2$\times$ extension.

\subsection{Extending LLM Context Window to 2048K}
\label{sec:extendto2048k}

\textbf{Progressive extension to 2048k}. In this section, we present an approach for expanding the context window of pre-trained LLMs from the conventional 4k tokens to beyond 2048k tokens. Our experiments show that with non-uniform positional interpolation, the context window can be enlarged by up to 8$\times$ without the need for fine-tuning. When a greater increase, such as 512$\times$, is desired, fine-tuning becomes essential. One approach involves identifying appropriate RoPE rescaling factors for the target context length (2048k) and subsequently performing fine-tuning. Nonetheless, this process entails significant computational costs, making it less feasible. Furthermore, in our experience, achieving effective fine-tuning for such substantial context window expansions is quite difficult (refer to Appendix for details).

Fortunately, LongRoPE demonstrates efficacy with both the base and fine-tuned extended LLM variants. To this end, we propose a streamlined, incremental approach that reaches the desired 2048k target length through only 1k fine-tuning steps, utilizing training sequences no longer than 256k.

$\Diamond$ \textit{LongRoPE-based search for scaling pre-trained LLM context to 256k}. Using LLaMA2 for illustration, we perform a search procedure targeting context window sizes of 128k and 256k. At this juncture, the context extension factors are 32$\times$ and 64$\times$, respectively.

$\Diamond$ \textit{Fine-tuning to 256k}. Following pre-training, we adapt the LLM to support a 256k context window. Initially, LLaMA2 undergoes 400 fine-tuning steps utilizing RoPE rescaling parameters set for 128k. After this, we substitute the RoPE rescale factors for those configured for 256k on the obtained checkpoint, and proceed with 600 further fine-tuning steps. This two-stage process demonstrates greater efficiency compared to a direct fine-tuning approach targeting 256k from the outset.

$\Diamond$ \textit{Expanding the fine-tuned extended LLM to 2048k using LongRoPE search}. In the concluding phase, we carry out an additional search based on the fine-tuned LLM set to a 256k context length. This process enables us to attain a substantially expanded context window of 2048k, all achieved without extra fine-tuning steps. The total extension factor reached is $512\times$.

\textbf{Restoring performance in shorter context windows}. Upon increasing the context window to a substantial 2048k length, we observe degraded results within the initial context window. This decline is an established challenge associated with positional interpolation~\cite{pi}, which constrains the position embeddings for the original context window to occupy a limited subspace in higher dimensions, leading to adverse effects on the language model's capabilities. When applying a 512$\times$ expansion factor, positions inside the standard 4k context window are densely compressed.

To address this issue, we conduct an additional evolution search on the expanded LLM, specifically tuning the RoPE rescale factors for shorter context spans such as 4k and 8k tokens. For these reduced context lengths, we lower the upper bound for the searched $\lambda$, since less positional interpolation is necessary. While running inference, the LLM adapts the associated RoPE rescale factors on-the-fly.

 \section{Experiments}

\label{sec:exp}
\subsection{Setup}
\textbf{Evaluation Tasks and models}. LongRoPE is integrated with LLaMA2-7B and Mistral-7B, and its effectiveness is assessed across three dimensions: (1) measuring the perplexity scores for language models when processing lengthy documents with extended context windows; (2) evaluating performance on the Passkey retrieval task, which gauges the models' proficiency in extracting a straightforward passkey amidst a large amount of irrelevant information; and (3) testing on conventional LLM benchmarks restricted to a context window of 4096 tokens.

\textbf{Fine-tuning}. For LLaMA2, the model is fine-tuned with a learning rate set to 2e-5 employing linear decay, using a global batch size of 32. The training process consists of 400 steps conducted on the Redpajama~\cite{redpajama} dataset, which is partitioned into 128k-token segments and surrounded by BOS and EOS tokens. Afterward, utilizing the resulting checkpoint, training continues for an additional 600 steps to extend the context window to 256k tokens. The training for the 128k context employs 8 A100 GPUs with the distributed training system~\cite{cube}, while the expanded 256k context requires 16 A100 GPUs. For Mistral, the setup involves maintaining a fixed learning rate of 1e-6 and a global batch size of 64. The procedures for both the 128k and 256k token models adhere to YaRN~\cite{yarn} settings, involving 400 steps on Together Computer's Long-Data Collections~\cite{mistraldata} with a 16k sequence length. Four A100 GPUs are used for this training.

\textbf{Search}. When the target window size is less than or equal to 256k, the parameters are set as follows: $P$=64, $N_1$=$N_2$=16, $p$=0.3, and $\mathcal{T}$=40. In each search iteration, the top 32 candidates are chosen for mutation and crossover operations. Perplexity evaluation employs five randomly selected samples from the PG19 validation set, ensuring that each sample meets the minimum target context length. For window sizes exceeding 512k, the population size, as well as the mutation and crossover candidate numbers, are reduced by half. In these cases, perplexity is computed using three randomly chosen validation samples from Pile-Books3~\cite{pile}.

\textbf{Baselines}. To achieve a 2048k context window, we employed fine-tuning strategies for models initialized with 128k and 256k context lengths. This results in LongRoPE-2048k (ft=128k) for LLaMA2 and LongRoPE-2048k (ft=256k) for Mistral. We assess these four models against leading methods for extending context windows, focusing on open-source LLMs that have undergone positional interpolation fine-tuning via PI, NTK, and YaRN techniques. The set of comparison models includes Together-32k~\cite{together}, Code LLaMA~\cite{codellama}, LongLoRA-full-FT-100k~\cite{longlora}, as well as YaRN-LLaMA and YaRN-Mistral~\cite{yarn}.

\subsection{Main Results}

\label{sec:mainresult}
\textbf{Language modeling of long sequences up to 256k tokens}. To initiate our evaluation, we benchmark against leading extended LLMs on sequence lengths up to 256k. Our assessment spans two benchmark datasets—Proof-pile~\cite{pg19} and the PG19~\cite{pile} test splits—to showcase generalization capability. We report perplexity across varying context lengths, utilizing a 256-token sliding window. For PG19, the complete test set comprising 100 documents is employed. In the case of Proof-pile, following the methodology in YaRN~\cite{yarn}, we randomly sample 10 instances, each no shorter than 128k tokens.

Tables~\ref{tbl:proof-pile} and~\ref{tbl:pg19} present a comparison of perplexity between LLaMA2 and Mistral models that have been extended using various interpolation techniques, evaluated on the Proof-pile and PG19 datasets, respectively. Two principal findings emerge: \textbf{(1)} Our extended models consistently exhibit lower perplexity scores as the evaluation length increases from 4k up to 256k, demonstrating their capability to effectively utilize longer contexts. \textbf{(2)} Remarkably, despite the substantial extension of the context window to 16$\times$ its original size—a scenario often detrimental to short-length performance—our LongRoPE-2048k models surpass leading baselines on context lengths up to 256k.

\textbf{Modeling languages on sequences over 2000k tokens}. To assess performance on exceptionally lengthy texts, we utilize the Books3~\cite{pile} dataset. For efficiency during evaluation, 20 books longer than 2048k tokens are chosen at random, and a 256k-token sliding window is applied.

Table~\ref{tbl:books3} demonstrates that LongRoPE effectively expands the context window of both LLaMA2-7B and Mistral-7B up to 2048k, while maintaining perplexity scores at shorter context lengths (8k-128k) that are on par with, or outperform, the baseline models. We note substantial variation in performance between the 2048k versions of LLaMA2 and Mistral. Specifically, Mistral exhibits superior results relative to baselines at smaller context windows, although its perplexity rises above 7 for lengths beyond 256k. Meanwhile, LLaMA2 shows expected behavior: perplexity systematically decreases with larger context sizes, aside from slight increases observed at 1024k and 2048k. For LLaMA2, LongRoPE-2048k yields stronger results when fine-tuned with a 256k window compared to 128k, which can be attributed to a lower secondary extension ratio (8$\times$ as opposed to 16$\times$). Conversely, Mistral benefits from fine-tuning at the 128k window size. This difference stems from the choice in Mistral’s 128k and 256k fine-tuning process: we adopted YaRN's methodology by employing a 16k training length, thereby influencing Mistral’s capacity to extend its context window following fine-tuning.

\textbf{Passkey retrieval}. Our investigation turns to the practical context window size relevant for generation tasks. We utilize a synthetic benchmark for passkey retrieval, originally introduced by~\cite{passkey}. Here, the objective is for the model to extract a randomly generated five-digit passkey embedded within an extended textual input. The precise prompt format is provided in the appendix. For evaluation, we conduct 10 separate runs of the task, positioning the passkey in a randomly selected spot, uniformly distributed throughout the available context length.

Figure~\ref{fig:passkey} presents a comparative analysis of retrieval accuracy across various baseline models. The retrieval performance of current approaches declines sharply, reaching zero for token lengths greater than 128k. However, LongRoPE-LLaMA2-2048k (ft=256k) demonstrates remarkable robustness, sustaining retrieval accuracy above 90\% throughout the range of 4k to 2048k tokens—even under the demanding scenario of extracting a passkey from over a million tokens. Similarly, LongRoPE-Mistral-2048k (ft=128k) achieves perfect retrieval accuracy (100\%) up to 1800k tokens, only decreasing to 60\% at 2048k, which is consistent with Table~\ref{tbl:books3}, where a minor rise in perplexity is observed at the 2048k length.

\textbf{Standard benchmarks under the original context window}. We assess the performance of LongRoPE-2048k models within their native context length by utilizing the Hugging Face Open LLM Leaderboard~\cite{huggingfaceboard}, considering both zero-shot and few-shot configurations. Our evaluation comprises 25-shot ARC-Challenge~\cite{allenai:arc}, 10-shot HellaSwag~\cite{zellers2019hellaswag}, 5-shot MMLU~\cite{mmlu}, and 0-shot TruthfulQA~\cite{lin2021truthfulqa}.

According to Table~\ref{tbl:commonsense}, our models attain results similar to those on the initial benchmark—which was intended for limited context length—and surpass the original Mistral's performance on TruthfulQA by a margin of +0.5\%. LongRoPE-LLaMA2-2048k, which underwent fine-tuning at 256k, exhibits somewhat greater decline in performance; nevertheless, its scores remain acceptable across the majority of evaluated tasks.

\subsection{Ablation Results}

\textbf{Assessing the second positional interpolation's efficacy.} Within the framework of our progressive extension methodology, a secondary non-uniform positional interpolation is executed via our search algorithm on the already fine-tuned and extended LLMs. To assess its utility, we perform evaluations on our LLaMA2-256k model after fine-tuning. This model is subsequently scaled to 512k, 1024k, and 2048k contexts employing PI and YaRN methods. According to the outcomes displayed in Table~\ref{tbl:secondsearch}, our non-uniform positional interpolation preserves perplexity levels, ensuring stability. Conversely, perplexity under PI and YaRN exhibits rapid escalation as the extension ratio increases.

\textbf{Enhancing performance at reduced context lengths.} In order to address the drop in effectiveness at shorter contexts, we fine-tune the RoPE scaling parameters for LongRoPE-2048k using our optimization procedure. This involves imposing stricter limits on the scaling factors considered in the search process, thereby reducing interpolation at 4k and 8k context lengths. Table~\ref{tbl:shortperformance} presents a side-by-side analysis of perplexity for LongRoPE-LLaMA2-2048k evaluated on the Proof-pile dataset at 4k and 8k lengths, as well as the mean accuracy across several LLM benchmarks. These findings reveal notable gains in model performance when operating with shorter context windows.

\textbf{Examination of the two types of non-uniformities}. In the final stage, we perform ablation studies to isolate the effects of each non-uniformity component on overall performance. Two experimental setups are employed: (i) scaling LLaMA2-7B to shorter sequence lengths of 16k and 32k utilizing different approaches—PI, optimizing only the RoPE dimension, and optimizing both non-uniformities; (ii) extending our fine-tuned LLaMA2 model, trained at 256k sequence length, up to 2048k, maintaining the same approach. Perplexity measurements are reported in the absence of further fine-tuning. The results presented in Table~\ref{tbl:searchdimension} demonstrate that adjusting non-uniformity within the RoPE dimension yields a substantial reduction in perplexity compared to linear interpolation as used in PI. Modifying non-uniformity in token position enhances performance at 16k and 32k lengths, but this advantage does not persist at the 2048k length—likely due to the challenges posed by such extended sequences. Simply keeping the initial tokens without interpolation becomes ineffective under these conditions; we propose to investigate this further in subsequent studies.

\section{Related Works}

Beyond techniques relying on position interpolation, this section reviews alternative lines of research. 

\textbf{Retrieval-based} methods leverage external memory modules to store extensive historical context, and employ retrieval mechanisms during inference to obtain relevant documents~\cite{longllama,wang2023augmenting,borgeaud2022improving}. These approaches often require substantial and explicit architectural changes to the underlying LLMs. In comparison, our approach remains more streamlined, introducing only minimal adjustments to positional embeddings. Furthermore, our method supports a broader range of long-context applications beyond retrieval, such as summarization of lengthy documents and few-shot learning.

\textbf{Attention-based context window extensions}. In addition to positional embedding interpolation, several studies have managed to enlarge the input context window without altering the original context length in LLMs by adjusting the attention mechanisms~\cite{han2023lminfinite, streamingllm,parallelwindow}. The main approach involves utilizing innovative attention masks to address the problem of attention explosion introduced by the inclusion of new positions. Notably, these strategies and methods based on positional interpolation are mutually supportive.

\textbf{Fine-tuning based methods} explore strategies for efficiently adapting pre-trained LLMs with altered position embeddings to handle extended context lengths. Previous research such as Code LLaMA~\cite{codellama}, LLaMA2 Long~\cite{xiong2023effective}, and ScaledRoPE~\cite{liu2023scaling} employ significantly larger base values for RoPE and perform fine-tuning at the desired sequence length. The approach introduced in this paper allows for adaptation across a broad range of target lengths, reaching up to lengths exceeding 2M tokens. In recent developments, as tuning for longer contexts (e.g., greater than 128k tokens) imposes heavy demands on GPU computation, solutions like LongLoRA~\cite{longlora} and PoSE~\cite{zhu2023pose} have been designed to alleviate these computational bottlenecks. Notably, our approach operates independently of these resource-efficient fine-tuning techniques.

\section{Conclusion}

This paper introduces LongRoPE, a technique designed to drastically increase the context window of LLMs to 2048k tokens—surpassing prior limits—while preserving their performance in original, shorter contexts. The approach leverages two distinct types of non-uniformities found within RoPE positional embeddings and utilizes a streamlined evolutionary search algorithm. This methodology yields two principal advantages: it supplies a robust foundation for subsequent fine-tuning and facilitates an 8$\times$ context window expansion even in the absence of fine-tuning. Furthermore, we develop a staged extension framework by fine-tuning LLMs with 256k-length context windows, ultimately scaling to 2048k without requiring additional fine-tuning. Comprehensive experimentation corroborates the efficacy of LongRoPE. We anticipate that LongRoPE-2048k models will unlock a variety of novel applications dependent on extended contexts and stimulate further exploration in this area.

\section*{Broader Impacts} The aim of this paper is to contribute to ongoing developments within Machine Learning. While our research may give rise to various societal effects, we do not consider any particular consequence necessary to emphasize in this section.

	\balance

\end{document}